每天早上你要给一条曲线交两个数：今天大概是多少，超过多少算异常。曲线可以是订单量、机房负载、气温、库存周转，也可以是模型上线之后的某个线上指标。手上有的就是这条曲线的历史，别的什么都没有。

本 Part 前面几章给了描述随机过程的词汇——轨道与截面、自协方差、平稳、遍历——但那些词汇回答的是``一个过程是什么样的''，不是``手上这条序列该怎么办''。从词汇到操作之间缺的那一段，就是本章要补的。

补法的核心可以用一句话概括：把序列一层层剥开，直到剩下的东西不可预测为止。

\begin{center}
\begin{tikzpicture}[x=1cm,y=1cm,font=\small,>=stealth]
  \draw[black!50,fill=blue!5] (0,0) rectangle (3.4,1.2);
  \node[font=\footnotesize,align=center,black!80] at (1.7,0.6) {观测序列\\\footnotesize 有趋势、有惯性};
  \draw[black!50,fill=blue!5] (6.0,0) rectangle (9.4,1.2);
  \node[font=\footnotesize,align=center,black!80] at (7.7,0.6) {平稳序列\\\footnotesize 还有自相关};
  \draw[black!50,fill=green!7] (12.0,0) rectangle (15.4,1.2);
  \node[font=\footnotesize,align=center,black!80] at (13.7,0.6) {白噪声\\\footnotesize 剩下的不可预测};
  \draw[->,black!65,thick] (3.5,0.6) -- (5.9,0.6);
  \draw[->,black!65,thick] (9.5,0.6) -- (11.9,0.6);
  \node[above,font=\footnotesize,black!75] at (4.7,0.72) {\((1-B)^d\)};
  \node[below,font=\footnotesize,black!60,align=center] at (4.7,0.46) {差分\\去趋势};
  \node[above,font=\footnotesize,black!75] at (10.7,0.72) {\(\theta(B)^{-1}\phi(B)\)};
  \node[below,font=\footnotesize,black!60,align=center] at (10.7,0.46) {扣掉自身历史\\与冲击回声};
  \node[font=\footnotesize,red!65,align=center] at (4.7,-0.6) {差过头会放大噪声};
  \node[font=\footnotesize,red!65,align=center] at (10.7,-0.6) {阶数过高会过参数化};
  \draw[->,black!45] (13.7,-0.08) -- (13.7,-1.35) -- (1.7,-1.35) -- (1.7,-0.08);
  \node[above,font=\footnotesize,black!60] at (7.7,-1.32) {残差不像白噪声，就退回来改设定};
\end{tikzpicture}

\emph{整章的骨架。两个算子各管一件事：\((1-B)^d\) 拿掉随每一次冲击永久累积的随机趋势，\(\theta(B)^{-1}\phi(B)\) 拿掉线性可预测的部分。走到最右端时，白噪声的方差就是这套模型能达到的预测误差下限；走不到最右端，说明还有结构没被建进去。两个红字提醒的是同一件事——每一步剥离都要付代价，剥多了比剥少了更难发现。}
\end{center}

沿着这条骨架，本章有五条判断需要立住。

第一条关于 AR。它做的事情不神秘，就是把序列对自己的历史做线性回归；有意思的是稳定性完全写在特征根里——根落在单位圆内，冲击几何衰减，系统有一个可回归的水平；根贴到单位圆上，冲击永不消失，序列变成随机游走。这与 Markov 链收敛速度由第二大特征值控制、线性递推的解由特征根决定、循环网络的长依赖由权重矩阵谱半径决定，是同一件事的不同外衣。

第二条关于 MA。它把冲击的记忆长度写死在 \(q\) 期，而 AR 的记忆无限长但按几何速率褪色。在可逆条件下两者可以互相改写，所以问``该用哪个''如果问的是表达力，答案是两者重合；真正的差别是参数个数。ARMA 的价值因此是简约性——用分子分母各几个参数逼近一整条无穷长的权重序列。

第三条关于差分，也是实践中最常被滥用的一步。差分能拿掉随机趋势，代价是它同时是个高通滤波器：低频被压成零，高频被放大一倍。差过头会凭空造出不可逆的 MA 结构、把噪声方差翻倍，而且诊断信号很微妙。更要紧的是趋势平稳与单位根平稳需要分开对待，方向错了不只是效率损失，是把信号本身差没了。

第四条关于前提。上面所有结论都建立在平稳性上，而真实序列常常不平稳——趋势、季节、结构突变、机制切换。好在不平稳是可观察的：自相关衰减得快不快，分段估计的均值方差是否可比，单位根检验怎么说（以及它在什么情况下说不准）。这条前提对做线上系统的人有个熟悉的名字，叫分布漂移。

第五条关于边界。非线性、状态切换、多变量强耦合、以及数据量足够时的跨序列学习模型，都是 ARMA 该让位的场合。这不是方法优劣，是可用信息量与模型容量之间的匹配问题。

五篇分别承担这些判断：

\begin{itemize}
\tightlist
\item \hyperref[sec:06-Random-Processes/06-ARMA-and-ARIMA/01]{``AR 模型用历史预测未来''} 用一张冲击响应图把 \(\phi\) 从拟合参数变成关于系统的断言，并把稳定性交给伴随矩阵的谱；
\item \hyperref[sec:06-Random-Processes/06-ARMA-and-ARIMA/02]{``MA 与 ARMA：冲击的持久影响''} 对照两条相关图给出定阶直觉，说明可逆性挑的是唯一写法，并用 Wold 分解交代 ARMA 在平稳世界里的位置；
\item \hyperref[sec:06-Random-Processes/06-ARMA-and-ARIMA/03]{``ARIMA：非平稳序列的差分处理''} 把差分的收益与代价一起算清楚，区分趋势平稳与单位根平稳，并说明单位根检验回答得了什么；
\item \hyperref[sec:06-Random-Processes/06-ARMA-and-ARIMA/04]{``ARMA 与其他方法的连接：什么时候用什么''} 是一篇选型判断，重点在于``可用信息量''与``偏离假设的程度''这两条轴；
\item \hyperref[sec:06-Random-Processes/06-ARMA-and-ARIMA/05]{``滞后回归与 ARDL：把回归变成动态模型''} 把外生变量接进来，用长期乘数和误差修正区分短期反应与最终水平，并说明为什么它仍然不是因果。
\end{itemize}

前三篇是主线，按顺序读收益最大；第四篇可以在真正要做选型决策时回头看；第五篇在手上不止一条序列、还握着价格或投放这类可控变量时才用得上。

如果只带走一样东西，建议是那张冲击响应图。它同时解释了平稳与非平稳的分界、差分为什么必要、过度差分为什么有害、以及长依赖为什么在循环网络里同样难处理。下一章离开时间序列，转向另一个方向：不再问``下一个值是多少''，而问``一条消息平均需要多少比特才能说清楚''——那是\hyperref[ch:053]{``信息论：从不确定性到可通信、可学习的边界''}的入口。
